\section{Conclusion}

This work builds on using model-predictions directly in hazard assessment of \citet{PosthumaImprovingHazard}, which replaces lognormal SSD curve-fitting with point predictions of LC50 for every species in the dataset. We extend that approach by quantifying the uncertainty attached to each prediction, making the predictions more trustworthy for practical use and giving a more faithful account of the variability in the testing data. The Hierarchical Bayesian Factorization Machine (BFM) provides, for every species--chemical pair, two distinct and interpretable uncertainty components: epistemic uncertainty, reflecting the model's lack of knowledge in under-explored regions of the chemical--species matrix, and aleatoric uncertainty, capturing the irreducible variability inherent to biological testing at the per-chemical level. This was achieved without loss of accuracy: the BFM reaches an RMSE of 0.70 log mg/L in grouped 5-fold cross-validation, improving on the 0.85 reported by \citet{PosthumaImprovingHazard}.

The two components behave as a trustworthy decomposition should. The hierarchical prior on the per-chemical precision parameter enables principled partial pooling: chemicals with abundant test data have their noise variance dominated by the empirical residuals, while data-poor chemicals are regularised toward the population mean. Epistemic uncertainty decreases with observation count, confirming that the model becomes more confident where more evidence exists, while aleatoric uncertainty shows no systematic trend with data, consistent with an irreducible noise term whose scale matches the spread of repeated laboratory tests. The combined intervals are well-calibrated at the coverage levels most relevant for regulatory use. Because the reducible component shrinks as data accrues while the irreducible component does not, the balance between the two depends on how well a chemical has been studied: predictions for sparsely tested chemicals are dominated by reducible model uncertainty, whereas predictions for well-studied chemicals are dominated by irreducible measurement noise.

We further show how this uncertainty propagates into the downstream use of ecotoxicity data for hazard assessment. Constructing SSDs directly from the ensemble of posterior samples yields SSDs with credible intervals that preserve the correlated epistemic structure across species, so the entire distribution shifts coherently with the shared parameters rather than individual points scattering independently. This enables uncertainty-aware derivation of Hazardous Concentrations such as the HC20, which correlate strongly with those from traditional lognormal curve-fitting while adding quantified credible intervals and covering all species in the dataset rather than only the tested subset. The width of these intervals communicates how much a derived hazard value should be trusted, and is widest for the data-poor chemicals where traditional small-sample SSDs are least reliable. Generating predictions with calibrated uncertainty for the 99.5\% of the chemical--species matrix that remains untested addresses a long-standing bottleneck in environmental risk assessment, equipping risk assessors with a distribution of plausible outcomes rather than a single point estimate, and supporting decisions that account for the evidence available for each chemical.

Two facts underlie this work: laboratory measurements are not perfect ground truth but noisy samples, and predictions from any statistical model carry uncertainty. Both can be modelled explicitly, and as shown here, doing so need not cost predictive accuracy. Neither the Bayesian reformulation that yields epistemic uncertainty nor the explicit per-chemical noise model that yields aleatoric uncertainty is specific to the factorization machine or to ecotoxicology. The same two steps can be applied to other models, and in other fields, wherever predictions are made from sparse, noisy, and unevenly sampled data and used to inform decisions, as recent uncertainty-aware work in adjacent areas has begun to do \citep{von2026uncertainty, iwasaki2026bayesian}. Treated this way, machine learning is not a black box: domain understanding and statistical modelling combine to show where a prediction can be trusted, where it cannot, and where the variability is an irreducible feature of the system that any assessment should take into account.

Finally, because the model reports where its predictions are least certain, it could guide the prioritisation of new tests towards the chemical--species combinations that would most reduce uncertainty. This is an uncertainty-guided strategy analogous to active learning, which could direct testing effort more efficiently and, over time, contribute to the goal of reducing animal testing.

\section{Limitations and future work}

Several limitations of this work should be noted, each indicating a direction for further development.

\begin{enumerate}
    \item Aleatoric uncertainty is modelled at the chemical level rather than the species or observation level. The precision parameter $\alpha_c$ is shared across all species and durations for a given chemical, so the model cannot represent the possibility that some species respond more variably to a given chemical than others. A more granular parameterisation could improve calibration but would increase the number of parameters and the data required to fit them. Resolving the aleatoric variance further, for example jointly across chemicals and species, is a natural next step.

    \item The model cannot extrapolate to entirely novel chemicals or species. The factorization machine relies on one-hot encoding of chemical and species identities, so predictions require both to have been observed at least once during training. Extending the model to untested chemicals would require richer molecular representations, such as molecular fingerprints or SMILES embeddings, together with phylogenetic features, and remains a major open challenge in the field.

    \item A mild over-prediction of aleatoric variance remains for data-poor chemicals. At low replicate counts the predicted standard deviation exceeds the empirical one by a pooled factor of roughly 1.2 (Table~\ref{tab:replicate_stats}), relaxing towards unity as replicate counts grow. As discussed in the results, this is largely explained by the downward bias of the empirical standard deviation when it is estimated from only a few replicates, rather than by miscalibration of the model. For regulatory use the residual effect is conservative, widening rather than narrowing the intervals on derived hazard concentrations.

    \item Only the LC50 mortality endpoint is considered. Incorporating additional endpoints, such as EC50 or NOEC, could improve predictive power by increasing the number of observations per chemical, but would require careful harmonisation of endpoint definitions.

    \item The model is predominantly data-driven, relying mainly on chemical and species identities together with a small set of features. This makes it sensitive to biases in the underlying dataset and limits its direct transfer to other settings. Applying the methodology to other public or proprietary datasets, and evaluating its applicability domain, would clarify how far the predictions and their uncertainties can be relied upon for new chemicals and species.
\end{enumerate}